% generated by analysis/paper_numbers.py -- do not edit
\newcommand{\teleportAuditCell}{SWAP ideal $-$0.071, hardware $+$0.040}
\newcommand{\teleportParagraph}{A teleported long-range CNOT on 104 qubits (logical depth 9) was compared with a SWAP ladder (depth 31) at distance 12. We replayed every SWAP-arm hardware generation exactly: beliefs rebuilt from the recorded frequencies, with the environment schedule matching the recorded one in every generation. The SWAP arm reads $+$0.040 on hardware where its ideal is $-$0.071; the sign is wrong in 72\% of 32 generations. The teleport arm reads $-$0.128$\pm$0.012 on hardware; its exact replay (mid-circuit measurement forces a per-shot simulation of 104 qubits) did not finish within our compute budget, so we report no ideal for it. The original write-up quoted an exact ideal of $-$0.29; its own script, rerun, gives a mean of $+$0.0042 on its reduced 18-qubit model with alternating sign. A depth comparison alone decides nothing; the exact replay does.}
\newcommand{\scanParagraph}{We measure the distance to L2 on two axes. \emph{Cost.} The most expensive exact reproduction of a hardware result in Table~\ref{tab:record} took 18~s, at most 146 core-seconds on eight cores. Taking $10^6$ core-years ($3.2\times10^{13}$ core-seconds) as a conservative threshold for ``out of classical reach'', the record sits 11.3 orders of magnitude below it (Fig.~\ref{fig:map}). For the genealogy family this distance does not shrink with better hardware, because $\chi$ stays 2 at every width. \emph{Depth.} The QCA is the only family with non-Clifford coupling on a lattice with loops. On the live 23-site region, belief propagation with $\chi\le16$ converges up to depth $t=3$ and fails from $t^*=4$. Hardware passed its pair witness at $t=2$, so the depth distance is 2 steps. Failure of belief propagation is not classical hardness: exact MPS reproduces the same 23 sites at $t=4$ in 68.53~s ($\chi=492$). Exact cost does rise steeply with width and depth (Table~\ref{tab:scan}): at 40 or more sites, $t=3$ costs 42 to 144~s ($\chi$ up to 482) and $t=4$ did not finish within 1200~s and 20~GB at any width. That is the direction in which an L2 candidate could appear. It needs a passing witness at the same width and depth, and today the witness at $t=2$ passed on one of four pairs.}
\newcommand{\ctParagraph}{The live run rolled twelve tumours (8 independent words and two sibling pairs) in one job on \texttt{ibm\_marrakesh}, with a CHSH test in the same job: $S=2.733\pm0.032$ (2048 shots per setting), 22.7$\sigma$ above the local bound. Every case is therefore P-L0. Each grown tumour starts from one Z snapshot of a 23-site QCA measured in a second job, whose pair witness passed marginally. We compared three levels with classical twins. \emph{Roll.} The independent size draws were (1, 0, 5, 2) for the bands $<6$, 6--8, 8--20 and $>20$~mm, against (4.4, 2.0, 1.2, 0.4) expected under the rulebook (Table~\ref{tab:roll}). Against 200{,}000 classical cohorts drawn from the same rulebook this gives $p=1.7\times10^{-4}$ (G statistic), and the joint log-likelihood falls at percentile 0.022\%. The 9-qubit roll compiles from 104 to 642 two-qubit gates (depth 1616); simulated Heron noise moves the ideal size mix (0.554, 0.247, 0.148, 0.051) to (0.438, 0.254, 0.194, 0.114), the same direction as the hardware. \emph{Seeds.} For each tile we compared the 48{,}000 hardware snapshots per arm with an ideal sampler of the same circuit, a noisy sampler and independent coins with the hardware marginals (Fig.~\ref{fig:twin}a). Hardware split halves agree (seed-size KS $D=0.010$, $p=0.18$); hardware differs from the ideal ($D=0.120$), noisy ($D=0.073$) and independent ($D=0.085$) twins, all $p\le9.5\times10^{-113}$. The mean pair-correlation RMS is 0.119 on hardware, 0.166 ideal, 0.114 noisy and 0.016 for independent coins. The seed structure comes from correlations that any classical sampler given the joint distribution reproduces, and that a marginals-only sampler loses. All 12 of 12 stored seeds were regenerated from the recorded hardware bits. \emph{Rendered nodules.} Four replicate cohorts per arm, each grown and rendered into the same Siemens-scanner host, gave 36 quantum-seeded and 36 twin-seeded nodules (Fig.~\ref{fig:twin}b). No feature differed: KS $p\ge0.88$, paired per-tile Wilcoxon $p\ge0.25$. Baseline diameter is equal by construction, because growth targets the rolled size. At these sample sizes the test detects $D\ge0.32$; smaller differences are not excluded.}
\newcommand{\ruleParagraph}{The case study separates what the quantum computer adds from what it costs. The Bell test certifies that the draw was not chosen. After rendering, the content cannot be told from its classical twin. Sampling the rulebook directly on hardware distorted the medical statistics ($p=1.7\times10^{-4}$). The rule that follows: certify the seed on the quantum computer, then sample the content classically from it, for example with a deterministic random bit generator keyed by the certified bits~\cite{barker2015}. The classical twin then \emph{is} the generator. The rulebook is reproduced up to sampling noise, and the job identifier, the $S$ value and the seed hash travel with every case as its replayable record. Quantum structure, where wanted (correlated sibling lesions, lineage-shaped seeds), is added only with a passing witness in the same job and labelled P-L1.}
\newcommand{\conclusionParagraph}{Today's quantum computers can do one thing for medical imaging that classical generators cannot: certify that synthetic validation cases were not chosen. In our live cohort the Bell value was $S=2.733$, 22.7$\sigma$ above the local bound. Every case is replayable from its seed, and its content cannot be told from a classical twin. The design rule is to certify the seed on the quantum computer and sample the content classically. From there the road to more complex quantum-native content has three further layers, each with a pass criterion: correlated families from one entangled lineage, compressed growth models, and seed structure beyond classical reach. The nearest open cell is the QCA at depth four and forty or more sites, which exact classical simulation on one machine did not finish. Quantum simulation of life itself is not on that road yet: every faithful circuit we ran was reproduced classically in at most 18~s. Certified, not superior, is the honest starting point, and it is already useful.}
\newcommand{\figOneCaption}{Every hardware result in our record, placed by its margin above the classical bound of its witness (vertical, symmetric log scale; marker = witness rung, colour = project) and by the measured cost of reproducing it exactly on one eight-core machine (horizontal, core-seconds). Dashed lines: cost tiers. Shaded: alive and beyond $10^6$ core-years, the region an L2 claim needs; it is empty. Lower strip: the QCA family in simulation only (exact MPS, ideal) by patch size and depth; arrows mark cells that did not finish within 1200~s and 20~GB.}
\newcommand{\figTwoCaption}{Classical-twin test of the procedurally generated CT nodules. (a)~Seed size over 48{,}000 QCA snapshots per arm: hardware, an ideal classical sampler of the same circuit, a noisy sampler (FakeMarrakesh noise model) and independent coins with the hardware marginals. (b)~Rendered nodules from four replicate cohorts per arm, quantum-seeded ($n=36$) vs twin-seeded ($n=36$), standardised per feature; two-sample KS $p$ per feature.}
